% concatenated from main.tex
\documentclass[12pt]{article}
\usepackage[margin=1.5in]{geometry} 
\usepackage{amsmath,amsthm,amssymb,amsfonts}
\usepackage{changepage}
\usepackage{tikz}
\usepackage[shortlabels]{enumitem}
\usetikzlibrary{matrix} 
 
\newcommand{\T}{\mathcal{T}}
\newcommand{\D}{\mathcal{D}}
\newcommand{\CC}{\mathcal{C}}
\newcommand{\oo}{\mathcal{O}}
\newcommand{\0}{\varnothing}
\newcommand{\Z}{\mathbb{Z}}
\newcommand{\N}{\mathbb{N}} 
\newcommand{\Q}{\mathbb{Q}} 
\newcommand{\R}{\mathbb{R}}
\newcommand{\F}{\mathcal{F}}
\newcommand{\p}{\mathfrak{p}}
\newcommand{\q}{\mathfrak{q}}
\newcommand{\m}{\mathfrak{m}}
\newcommand{\ann}{\textnormal{ann}}
\newcommand{\hm}{\textnormal{Hom}}
\newcommand{\id}{\textnormal{id}}
\newcommand{\ind}{\textnormal{Ind}}
\newcommand{\im}{\textnormal{im }}
\newcommand{\asdim}{\textnormal{asdim}}
\newcommand{\asi}{\textnormal{asi}}
\newcommand{\inv}{^{-1}}
\newcommand{\st}{\ | \ }
\newcommand{\tor}{\textnormal{Tor}}
\newcommand{\len}{\mathcal{l}}
\newcommand{\gal}{\textnormal{Gal}}
\newcommand{\leg}[2]{\left( \frac{#1}{#2} \right)_n}
\newcommand{\sleg}[2]{\left( #1/#2 \right)_n}
\newcommand{\pp}{(\oo/\p)^\times}
\newcommand{\qq}{(\oo/\q)^\times}
\newcommand{\tr}{\textnormal{Tr}}
\newcommand{\cay}{\textnormal{Cay}}
\newcommand{\PMEU}{PME\textsubscript{U}}
\newcommand{\res}{\upharpoonright}

\newtheorem{theorem}{Theorem}
\newenvironment{problem}[2][\!\!\!]{\begin{trivlist}
\item[\hskip \labelsep {\bfseries #1}\hskip \labelsep {\bfseries #2.}]}{\end{trivlist}}
%If you want to title your bold things something different just make another thing exactly like this but replace "problem" with the name of the thing you want, like theorem or lemma or whatever

\newtheorem{corollary}{Corollary}
\newtheorem{lemma}{Lemma} 
\newtheorem{prop}{Proposition}  
\newtheorem{prob}{Problem} 
 
\title{Disjoint Borel annuli that every injective ray meets infinitely often at Borel asymptotic separation index one}
\author{Felix Weilacher}
\date{arxiv:2104.14646v2; CC BY 4.0}
\begin{document}
\maketitle
\noindent\textit{Editorial note.} The counted unit is the article's own Lemma lem:rainbow (source lines 173-175) together with the proof that begins on the next line (177-186); the pairing is the ordinary adjacent one.\par\medskip
\noindent\textit{Editorial note.} Nothing outside the retained definitions at source lines 123-129 and 169 is used by this proof, so the package carries no support lemma and no bibliography.\par\medskip
\noindent\textit{Editorial note.} The article numbers this as its first lemma and the retained declaration block reproduces that numbering; the original label lem:rainbow is carried in the span metadata.\par\medskip
\noindent\textit{Editorial note.} No retained line contains a citation, a figure or an included file.\par\medskip
\noindent\textit{Editorial scope.} One original result: the article's Lemma lem:rainbow with its complete author proof, plus the article's own definitions of the Borel asymptotic separation index and of the ball and annulus notation used in that proof. The same section's Lemma lem:Z, the corollaries to Theorem th:d+1 and the two already counted claims from this source are not counted here and are not used as premises.\par\medskip
\noindent\textit{Allowed standard prerequisites (editorial).}\begin{itemize}
\item Locally finite Borel graphs and their distance graphs. A locally finite Borel graph on a standard Borel space, its distance N-graph G\textasciicircum{}\{<=N\}, Borel partitions, and connected components of an induced subgraph; all of this is retained from the article's own Section 2 passage at source lines 123-129.. The proof builds the annuli inside a Borel partition and reasons about G\textasciicircum{}\{<=N\}-components.
\item Borel asymptotic separation index one. asi\_B(G)=1 means that for every N the space splits into two Borel pieces U\_0, U\_1 with G\textasciicircum{}\{<=N\} restricted to each piece having only finite connected components (source lines 125-127). No consequence of the invariant beyond this witnessing partition is assumed.. The scale N=4d+1 and the partition of X into U\_0 and U\_1 at source line 180 come from this definition.
\item Elementary path counting in a locally finite graph. Path distance to a fixed set changes by at most one along consecutive vertices of an edge, and in a locally finite graph only finitely many vertices lie within distance 2d of a finite set.. Used at source lines 182 and 184 to force a ray out of the ball B(U\_0,2d) and then to trap it in each annulus.
\end{itemize}\medskip
We start with the discrete versions of these notions: Let $G$ be a locally finite graph on a set $X$, and $N$ a positive integer. Let $G^{\leq N}$ denote the \textit{distance $N$-graph} of $G$. This is the graph which sets two vertices adjacent if and only if there is a path from one to the other of length at most $N$.

The \textit{asymptotic separation index} of $G$, denoted $\asi(G)$ is the infimum over $s \in \omega$ such that the following holds: For every positive integer $N$, $X$ can be partitioned into $s+1$ sets, say $U_0,\ldots,U_s$, such that for each $0 \leq i \leq s$, the induced graph $G^{\leq N} \res U_i$ has only finite connected components. The \textit{asymptotic dimension} of $G$, denoted $\asdim(G)$, is defined similarly, except now the connected components of each $G^{\leq N} \res U_i$ are required to have a uniform (finite) bound on their diameter. 

 Let $G$ now be a Borel graph on a standard Borel space $X$. The \textit{Borel asymptotic separation index} and \textit{Borel asymptotic dimension} of $G$, denoted $\asi_B(G)$ and $\asdim_B(G)$ respectively, are defined exactly as their discrete counterparts were, except that now all the $U_i$'s are required to be Borel sets.
 
 Now let $\Gamma$ be a finitely generated group, and $a:\Gamma \curvearrowright X$ a free Borel action. It is easy to verify that none of the above numbers for the graph $G(a,S)$ depend on the choice of generating set $S$ for $\Gamma$. We can thus talk about the (Borel) asymptotic separation index and (Borel) asymptotic dimension of the action $a$ itself.

Let $G$ be a graph on a set $X$. An \textit{injective $G$-ray} is an injective infinite sequence $x_0,x_1,\ldots \in X$ such that $(x_i,x_{i+1}) \in G$ for all $i \in \omega$. Given a set $U \subset X$ and an $n \in \omega$, let us write $B(U,n)$ for the ball of radius $n$ around $U$. That is, the set of all points $x \in X$ whose path distance to $U$ is less than or equal to $n$. For $m \in \omega$, let us write $A(U,n,m) = B(U,m)-B(U,n)$. One can check that if $G$ and $U$ are Borel, these sets are always Borel.

\begin{lemma}\label{lem:rainbow}
Suppose $G$ is a locally finite Borel graph on a space $X$ with $\asi_B(G) = 1$, and $d \in \omega$. Then we can find pairwise disjoint Borel sets $V_1,\ldots,V_d \subset X$ such that any injective $G$-ray contains infinitely many edges within each $V_i$.
\end{lemma}

\begin{proof}
Since the tail of an injective $G$-ray is still an injective $G$-ray, it suffices to prove this with ``infinitely many edges'' replaced by ``an edge''. 

Let $X = U_0 \sqcup U_1$ be a Borel partition of $X$ witnessing $\asi_B(G) = 1$ for $N = 4d+1$. That is, such that the graphs $G^{\leq N} \res U_0$ and $G^{\leq N} \res U_1$ both have only finite connected components. Set $V_i = A(U_0,2(i-1),2i)$ for each $1 \leq i \leq d$. Also set $W = X \setminus B(U_0,2d) \subset U_1$. As was mentioned above, these are all still Borel. Let us now show these $V_i$'s work.

Let $x = (x_n)$ be an injective $G$-ray. We first show that $x$ contains some point from $W$. Suppose not. Then for each $n$, there is a $y_n \in U_0$ whose path distance to $x_n$ is less than or equal to $2d$. Now for each $n$, going from $y_n$ to $x_n$, then to $x_{n+1}$, then to $y_{n+1}$, produces a path from $y_n$ to $y_{n+1}$ of length at most $2d+1+2d=N$, so all the $y_n$'s are in the same $G^{\leq N} \res U_0$-connected component. By definition of $U_0$, then, the set $\{y_n \mid n \in \omega\}$ is finite. But now $x$ gives infinitely many points within distance $2d$ of this set, contradicting local finiteness.

Since $W \subset U_1$, the above shows in particular that $x$ contains some point from $U_1$. Similarly, $x$ must contain some point from $U_0$. Let $x_n \in W$ and $x_m \in U_0$. We can assume $m < n$. Now the path distance from $x_i$ to $U_0$ changes by at most 1 each time $i$ is incremented by 1. Thus, since $x_m$ has distance 0 from $U_0$ and $x_n$ has distance greater than $2d$ from $U_0$, for each $1 \leq i \leq d$, there must be some $m < j < n$ such that $x_{j}$ and $x_{j+1}$ have distance $2i-1$ and $2i$ respectively from $U_0$. Then $(x_j,x_{j+1})$ is an edge of $x$ contained in $V_i$, as desired.

\end{proof}

\end{document}
